% =====================================================================
\chapter{Autoevaluación}
\label{ap:autoevaluacion}
% =====================================================================

\section{Cálculo de la nota}

La nota es la suma de los puntos de la parte~A y de la parte~B, sobre 100. El examen se aprueba con \textbf{70 puntos} (70\,\%). Aunque se apruebe, si en alguna parte o en algún módulo se obtiene menos del 50\,\% de los puntos posibles, debe repasarse ese contenido antes de continuar.

\begin{table}[ht]
\centering\small
\caption{Puntuación por parte.}\label{tab:partes}
\begin{tabular}{@{}lllcc@{}}
\toprule
Parte & Contenido & Módulos & Puntos & Obtenidos \\
\midrule
A & Cuestiones (20 $\times$ 1~p) & M12--M15 & 20 & \\
B.1 & Bloque decoder & M14 & 15 & \\
B.2 & Atención con máscara causal & M13 & 15 & \\
B.3 & BPE a mano & M12 & 12 & \\
B.4 & Parámetros, FLOPs y memoria & M14 & 15 & \\
B.5 & RoPE y coste de la atención & M13, M14 & 11 & \\
B.6 & Entrenamiento del mini-GPT & M15 & 12 & \\
\midrule
\textbf{Total} & & & \textbf{100} & \\
\bottomrule
\end{tabular}
\end{table}

\begin{table}[ht]
\centering\small
\caption{Puntuación por módulo. En el ejercicio B.5, el apartado~(a) corresponde a M14 y los apartados (b) y~(c), a M13.}\label{tab:modulos}
\begin{tabular}{@{}llccc@{}}
\toprule
Módulo & Elementos & Puntos & 50\,\% & Obtenidos \\
\midrule
M12. Tokenización & A.1--A.5; B.3 & $5+12=17$ & $8{,}5$ & \\
M13. Atención & A.6--A.10; B.2; B.5(b,c) & $5+15+6=26$ & $13$ & \\
M14. Transformer & A.11--A.18; B.1; B.4; B.5(a) & $8+15+15+5=43$ & $21{,}5$ & \\
M15. Mini-GPT & A.19--A.20; B.6 & $2+12=14$ & $7$ & \\
\midrule
\textbf{Total} & & \textbf{100} & & \\
\bottomrule
\end{tabular}
\end{table}

\section{Puerta de salida de la fase III}

Para pasar a la fase~IV deben cumplirse las tres condiciones del método de estudio del curso:
\begin{enumerate}[label=\arabic*.]
  \item en el test de cada módulo (M12, M13, M14 y M15), al menos 9 aciertos de 12;
  \item en la autoevaluación del laboratorio interactivo de la fase, al menos un 70\,\%;
  \item en este examen de hito, al menos 70 puntos de 100.
\end{enumerate}
El proyecto de fase se considera superado con al menos 70 puntos en la rúbrica del \cref{tab:rubrica}.

\section{Qué hacer si no se alcanza el umbral}

\begin{itemize}[leftmargin=1.4em,label=\textendash]
  \item Localizar en el \cref{tab:repaso} las secciones correspondientes a cada fallo y anotar el error en el registro de errores, clasificado como de concepto, de cálculo o de despiste.
  \item Releer esas secciones y rehacer, sin consultar la solución, los ejercicios $\star\star$ y $\star\star\star$ del tema correspondiente de los apuntes del módulo.
  \item Repetir el test del módulo afectado hasta alcanzar 9 de 12 y volver a resolver, al cabo de unos días, los ejercicios de la parte~B en los que se perdió más de la mitad de los puntos.
  \item Si la nota está por debajo de 50, repasar la fase completa siguiendo el ciclo del método de estudio (teoría, laboratorio, práctica, ejercicios, test y tarjetas) antes de repetir el examen.
\end{itemize}

\begin{table}[ht]
\centering\small
\caption{Secciones que deben repasarse según los fallos.}\label{tab:repaso}
\begin{tabularx}{\textwidth}{@{}>{\raggedright\arraybackslash}p{4.3cm}X@{}}
\toprule
Fallos & Repasar \\
\midrule
A.1, A.4; B.3(b,c) & M12, §2.1 (entrenamiento de BPE) y §2.2 (codificación) \\
A.2, A.3; B.3(d) & M12, §1.3 (UTF-8), §1.5 (consecuencias para el español) y §2.4 (BPE a nivel de byte) \\
A.5 & M12, §5.5 (secuencialidad y paso al Transformer) \\
A.6; B.2(a,b,d) & M13, §2.1 (forma matricial de la atención) \\
A.7; B.2(f) & M13, §2.3 (por qué se divide por $\sqrt{d_k}$) \\
A.8; B.2(c,e) & M13, §3.1 (máscaras aditivas) y §3.2 (causalidad) \\
A.9; B.5(b) & M13, §4.2 (parámetros de la atención multi-cabeza) \\
A.10; B.5(c) & M13, §4.3 (coste en operaciones) y §4.4 (memoria y contexto largo) \\
B.1(a--d) & M14, §1.2 (bloque decoder-only), §2.1 (flujo residual) y §3.1 (MLP posición a posición) \\
A.11, A.12; B.1(e) & M14, §2.2 (pre-norm frente a post-norm) y §2.3 (crecimiento del flujo) \\
A.13; B.4(a) & M14, §3.3 (SwiGLU) \\
A.14, A.15; B.5(a) & M14, §4.3 (RoPE) \\
A.16 & M14, §1.4 (embeddings atados) y §5.2 (parámetros del modelo completo) \\
A.17, A.18; B.4(b--e) & M14, §5.1 a §5.3 (parámetros y memoria) y §6.2 y §6.3 (FLOPs y regla $6ND$) \\
A.19; B.6(a,b) & M15, §1.4 (pérdida de la predicción uniforme) y §2.3 (recuento de parámetros) \\
A.20; B.6(c,d) & M15, §1.3 (perplejidad), §4.2 (puntos de control) y §4.4 (lectura de curvas) \\
B.6(e) & M15, §5.2 (temperatura) \\
\bottomrule
\end{tabularx}
\end{table}

\paragraph{Siguiente paso.} Superada la puerta de salida, se continúa con la fase~IV, módulo~16: preentrenamiento a escala.
